\chapter{Implementation and Training Procedure}

\section{Software Organization}
The implementation is divided into reusable modules:
\begin{itemize}
\item \texttt{src/amc/data.py} loads RadioML data, normalizes I/Q samples,
      and creates stratified splits.
\item \texttt{src/amc/model.py} defines the Phase 1 CNN baseline.
\item \texttt{src/amc/communications.py} implements modulation, channel,
      constellation estimation, and demodulation.
\item \texttt{scripts/train.py} trains and saves the Phase 1 CNN checkpoint.
\item \texttt{scripts/evaluate.py} evaluates the Phase 1 classifier.
\item \texttt{scripts/run\_phase2\_link.py} executes the complete Phase 2
      transmitter--channel--receiver experiment.
\item \texttt{scripts/build\_parameter\_figures.py} generates parameter,
      architecture, and training figures.
\end{itemize}

\section{Phase 1 Dataset Preparation}
RadioML 2016.10a contains 11 modulation classes and SNR values from -20 dB
to +18 dB. Every example contains 128 complex samples represented by two
channels. Per-example power normalization is applied before training:
\begin{equation}
\tilde{\mathbf{x}} =
\frac{\mathbf{x}}
{\sqrt{\operatorname{mean}(\mathbf{x}^2)+\epsilon}}.
\end{equation}
The dataset is divided into 70\% training, 15\% validation, and 15\% testing
using deterministic stratification.

\section{Phase 1 CNN}
The baseline receives a $(2,128)$ tensor. It contains three convolutional
blocks with 64, 128, and 256 filters, followed by adaptive average pooling
and a dropout-regularized fully connected layer. The total number of
trainable parameters is 144,331. Batch normalization and ReLU activations
improve optimization, while pooling reduces the temporal representation.

\section{Training Configuration}
\begin{table}[H]
\centering
\caption{Phase 1 CNN training configuration.}
\begin{tabular}{p{6cm} p{7cm}}
\hline
\textbf{Parameter} & \textbf{Value} \\
\hline
Dataset & RadioML 2016.10a \\
Input & 2 channels $\times$ 128 samples \\
Number of classes & 11 \\
Data split & 70\% / 15\% / 15\% \\
Normalization & Per-example power normalization \\
Loss & Class-weighted cross-entropy \\
Optimizer & AdamW \\
Learning rate & $3\times10^{-4}$ \\
Weight decay & $10^{-4}$ \\
Scheduler & Cosine annealing \\
Batch size & 1024 \\
Epochs & 8 in the reported run \\
Random seed & 42 \\
Trainable parameters & 144,331 \\
\hline
\end{tabular}
\end{table}

\section{Phase 2 Link Configuration}
\begin{table}[H]
\centering
\caption{Phase 2 communication-link configuration.}
\begin{tabular}{p{6cm} p{7cm}}
\hline
\textbf{Parameter} & \textbf{Value} \\
\hline
Modulations & BPSK, QPSK, 8-PSK, 16-QAM \\
Frame length & 256 symbols \\
Channel & Block Rayleigh fading and AWGN \\
Impairments & Phase offset and frequency offset \\
Synchronization & Pilot-assisted simulation estimate \\
SNR sweep & -4 to +16 dB, 2 dB increment \\
Frames per SNR & 200 \\
Metrics & BER, SER, packet success, throughput, decision accuracy \\
Random seed & 42 \\
\hline
\end{tabular}
\end{table}

\section{Training and Evaluation Workflow}
The Phase 1 model is trained independently and evaluated on the held-out
RadioML test set. Phase 2 uses a controlled communication simulation so that
the effect of modulation selection can be measured without mixing it with
uncontrolled dataset changes. Each frame has a known transmitted bit
sequence, which enables exact BER and packet-success calculation.

\section{Verification}
The repository includes four automated tests. They verify reproducible data
splitting, input-length validation, modulation round trips, and clean
constellation identification. Python compilation and whitespace checks are
also executed before the final results are generated.
